\section{Introduction}

The GenAI4World workshop asks how generative AI systems should be evaluated for
global tasks, contexts, vulnerabilities, and usability beyond standard accuracy
metrics \citep{genai4world2026cfp}. This question matters because many real
users do not interact with LLMs through clean monolingual benchmark items. A
single request may contain an instruction in one language, content in another,
an English file name embedded in Hindi written in Latin script, or a local date,
currency, or honorific convention. A model may understand the literal task and
still violate the pragmatic contract that makes the response usable.

Consider the following Spanish-English editing request:
\begin{quote}\small
\texttt{Me ayudas a mejorar este texto para que suene natural? ``I can't join
the call today, but I can share my update by email.''}
\end{quote}
A globally aligned assistant should infer that Spanish is the instruction
language and that the quoted English sentence is the object of editing. Unless
translation is requested, the edited sentence should remain English. A response
that translates the sentence into Spanish is fluent and superficially helpful,
but it is unusable for the user's intended workflow. The user must diagnose the
failure, issue a corrective follow-up such as ``Please keep the rewritten text in
English,'' and check whether the repair succeeded. We call this extra burden the
\emph{Re-prompt Tax}.

Existing multilingual evaluations have greatly expanded language coverage and
have shown that translated benchmarks can be culturally distorted
\citep{singh2024globalmmlu,xuan2025mmluprox,huang2025benchmax}. Other work has
studied multilingual instruction following \citep{zhou2023ifeval,li2025xifbench}
and output-language alignment in code-switched interactions \citep{oh2026ola}.
These efforts are essential, but they leave an interactional gap. A user-facing
system is not only a question-answering engine; it is a collaborator in a
multi-turn workflow. When a first response violates a language, script,
preservation, register, or locale constraint, the cost is paid by the user in
extra turns, extra tokens, and uncertainty about whether the model will repair
the failure.

This paper treats multilingual usability as an \emph{interaction contract}
between a user and a model. The contract may specify, explicitly or implicitly:
(1) the task, (2) the expected response language, (3) the language of content to
be edited or summarized, (4) script or romanization requirements, (5) literal
spans that must not be changed, and (6) register or locale constraints. A model
can violate this contract even when a semantic task label such as ``rewrite'' or
``summarize'' is satisfied. Measuring only one-turn task success therefore
understates the burden imposed on multilingual and code-switched users.

Our contributions are:
\begin{itemize}
    \item We define \textbf{Re-prompt Tax}, a metric family for hidden
    multilingual repair burden: first-turn global alignment, re-prompt turn tax,
    Repair@1/2, unresolved rate, and repair-normalized token tax.
    \item We construct \textsc{RePromptTax-Stress-v0.2}, a balanced 120-item
    synthetic stress pilot over three language pairs and four task families:
    implicit editing preservation, output-language inference, quote
    preservation, and script/register/locale constraints.
    \item We evaluate three GPT-4.1-family API models \citep{openai2025gpt41}
    and full 120-item current-model refreshes for \texttt{gpt-5.4-mini} and
    \texttt{gpt-5.5}; under \texttt{gpt-5.5}, the contract improves first-turn
    global alignment from 81.7\% to 98.3\% while leaving inspectable residual
    failures.
    \item We provide diagnostics beyond the headline table: language slices,
    family-level failure modes, paired sign-test sensitivity, token-burden
    analysis, prompt-mechanism and repair-realism diagnostics, item-consistency
    analysis, paired judge refreshes, and launch-ready human/native review
    protocols with conservative claim gates.
\end{itemize}

We intentionally keep the claim boundary narrow. The pilot is synthetic and
stress-oriented, not a prevalence estimate for all global users. Its purpose is
to make a neglected class of repair costs measurable and reproducible.
